% Ch. 14 : VM et conteneurs, piles comparées
\documentclass[border=6pt]{standalone}
\usepackage{coursfig}
\begin{document}
\begin{tikzpicture}[x=1mm,y=1mm]
  % pile VM
  \node[layer=Ocre, minimum width=80mm, anchor=south west] at (0,0) {Matériel};
  \node[keybox=Enc, minimum width=80mm, anchor=south west, rounded corners=0pt] at (0,10) {Hyperviseur};
  \foreach \x/\k in {2/1, 42/2}{
    \node[layer=Brick, minimum width=36mm, minimum height=15mm, anchor=south west, align=center] (g\k) at (\x,21) {OS invité \k\\[-.3mm]{\normalsize\color{Muted}noyau + userland}};
    \node[layer=Plum, minimum width=36mm, anchor=south west] at (\x,36) {App \k};
  }
  \node[ftitle, anchor=south] at (40,59) {MACHINES VIRTUELLES};
  % pile conteneurs
  \begin{scope}[xshift=100mm]
    \node[layer=Ocre, minimum width=80mm, anchor=south west] at (0,0) {Matériel};
    \node[layer=Brick, minimum width=80mm, anchor=south west] at (0,10) {OS hôte : \textbf{un seul noyau partagé}};
    \node[keybox=Enc, minimum width=80mm, anchor=south west, rounded corners=0pt] at (0,20) {Moteur de conteneurs};
    \foreach \x/\k in {2/1, 42/2}{
      \node[layer=Sea, minimum width=36mm, minimum height=15mm, anchor=south west, align=center] at (\x,31) {Conteneur \k\\[-.3mm]{\normalsize\color{Muted}userland seul}};
      \node[layer=Plum, minimum width=36mm, anchor=south west] at (\x,46) {App \k};
    }
    \node[ftitle, anchor=south] at (40,59) {CONTENEURS};
  \end{scope}
  \node[note, anchor=north, text width=76mm, align=flush center] at (40,-4) {chaque VM paie un noyau complet : isolation forte, coût élevé};
  \node[note, anchor=north, text width=76mm, align=flush center] at (140,-4) {le noyau est partagé : léger, isolation moindre};
\end{tikzpicture}
\end{document}
